% OpenRDMA Evaluation：英文论文正文，可由主论文直接 \input。
% 已有数值的唯一来源是 reference/evaluation-preview.pdf；不要把 planned 表格当作结果。
\providecommand{\EvalFigureRoot}{report/figures}
\providecommand{\EvalPlanned}[1]{}

\section{Evaluation}
\label{sec:eval:overview}

We evaluate the implemented datapath at the boundary that the current evidence
supports: cycle-accurate RTL execution with abstract DMA interfaces, coupled to
a packet-level network model.  Five questions guide the evaluation.  Where does
the TX pipeline spend time?  When does packet spraying help?  How much path skew
can the receiver absorb?  What does ordered completion cost?  Finally, what
implementation cost has been established, and what remains to be measured on a
target FPGA?  The results expose both benefits and limits.  In the tested
single-QP workload, packet spraying raises network goodput by up to $2.48\times$;
with four active QPs and 4-KiB messages, however, spraying delivers 26.4\%
lower goodput than fixed placement.  The trace combines byte imbalance with a
different packet-emission schedule, so the experiment does not isolate their
individual contributions.  Receive-window pressure and global retirement are
separate constraints that additional path bandwidth alone does not remove.

\subsection{Methodology and Evidence Boundary}
\label{sec:eval:method}

\textbf{Archived co-simulation.}
The measured batch connects a Verilator simulation of the TX RTL, an ns-3
network, and a Verilator simulation of the RX RTL through Unix sockets
\cite{verilator,ns3}.  TX exports each frame with the simulated time at which
its final beat leaves the packet scheduler.  ns-3 injects the frame at that
time, models serialization, queuing, and propagation, and forwards the arrival
event to RX.  Timestamped advance messages coordinate the three processes.
Thus, the network may reorder frames without replacing the production times
chosen by the RTL.  Table~\ref{tab:eval:configuration} summarizes the archived
configuration.

\begin{table}[H]
  \centering
  \caption{Archived RTL/network co-simulation configuration.  Link rates and
  the 2-ns RTL cycle are model parameters, not measurements of a placed FPGA.}
  \label{tab:eval:configuration}
  \small
  \begin{tabular}{p{0.15\textwidth}p{0.79\textwidth}}
    \toprule
    Component & Configuration \\
    \midrule
    Toolchain & BSC 2023.01; Verilator 5.048; ns-3 3.48; cocotb 2.1.0.dev0 \\
    RTL & 2-ns model cycle; 4 QPs; 4 receive lanes; messages up to 16 KiB \\
    Receive state & 16-packet window per QP; 16 payload slots per lane;
      4 Ordered Queue entries per QP \\
    Network & 2 nodes; 4 independent 100-Gb/s point-to-point links;
      100-packet queue per link \\
    Path delay & 200 ns on all paths for rate/policy comparisons;
      controlled skew for receive-pressure experiments \\
    DMA model & TX completions returned in reverse batches of up to 8 requests;
      RX completion acknowledged after payload collection \\
    Packet profile & Private Ethernet/IPv4/UDP/BTH/RETH; ICRC and FCS omitted \\
    \bottomrule
  \end{tabular}
\end{table}

In this profile, PMTU is a packetization budget rather than pure payload size:
the first packet can carry PMTU$-28$ B and a continuation packet PMTU$-12$ B;
Ethernet/IPv4/UDP framing is outside this budget.  For example, a 4-KiB message
at a 4-KiB PMTU becomes 4,068 B plus a 28-B tail.
The model exercises request processing, abstract DMA, packetization, path
selection, receive-window handling, placement, and ordered retirement.  It does
not include a physical PCIe/GPU path, a switched congestion topology, ICRC/FCS,
or an ACK-to-TX retransmission loop.  The four equal-rate links isolate path
selection; they do not evaluate SRv6 encapsulation.  This scope is narrower than
deployment-scale MRC studies~\cite{mrc-srv6-paper}.

\textbf{Metrics.}
TX latency begins when RTL accepts a request and ends at either the first frame's
start of frame (SOF) or the latest end of frame (EOF) among all packets of that
message.  TX effective payload rate divides batch payload bits by the interval
from the first accepted request to the last emitted frame.  Network goodput
instead spans first frame injection to final frame arrival.  Per-message network
completion spans that message's first injection to its final arrival, excluding
request processing before injection and RX processing after arrival.  Because
these metrics use different boundaries, their numeric rates must not be
subtracted as ``network overhead.''

All plotted points are deterministic finite batches without a separate warm-up
interval; they are not steady-state NIC throughput.  The eight-message policy
comparison supports paired means, not tail-latency estimates.  Coordination was
checked with 100-ns, 1-$\mu$s, and 5-$\mu$s advance windows, which required 12,
3, and 2 rounds but produced identical injection times, paths, and arrivals.
Free-clock execution and offline replay produced the same events, and 20
per-path packet captures matched reference bytes and timestamps.

\begin{table}[tbp]
  \centering
  \caption{Functional checks associated with the archived experiment batch.
  RX short and long suites overlap and are not additive coverage.}
  \label{tab:eval:validation}
  \scriptsize
  \begin{tabular}{@{}p{0.27\columnwidth}p{0.43\columnwidth}p{0.20\columnwidth}@{}}
    \toprule
    Check & Exercised behavior & Result \\
    \midrule
    TX regression & Request, DMA, packets, scheduler & 64 pass \\
    RX short & Window, placement, retirement, Flag & 37 pass; 1 skip \\
    RX long & Extended receive/completion & 38 pass \\
    Online reference & 6 messages, 4 QPs, unequal paths & 11/11 frames; 0 drops \\
    Timing consistency & 3 windows, free-clock, replay & Exact match \\
    Reverse DMA & Complete 3,2,1,0 across QPs & Retire 0,1,2,3 \\
    Queue capacity & 1, 2, 4 Flags; fifth blocked & Overflow caught \\
    \bottomrule
  \end{tabular}
\end{table}

The one short-suite skip in Table~\ref{tab:eval:validation} is the M6
long-only reorder case; it runs in the long suite.  The active repository has
since evolved into M0--M8 TX/RX units plus the M9--M11 software-coordinated
retransmission seam.  The latter is not exercised by this batch.  Moreover, the
current tree preserves the RTL and cocotb suites but not the external ns-3
adapter, raw event archive, or result parser used here.  We therefore treat the
figures as archived evidence rather than claim push-button reproduction from
the current checkout.

\EvalPlanned{
\begin{table}[tbp]
  \centering
  \caption{\textbf{P-XMODEL --- PLANNED MEASUREMENT -- NO DATA.}
  Cross-model trace closure required before replacing the archived batch.}
  \label{tab:eval:planned-xmodel}
  \footnotesize
  \begin{tabular}{llllll}
    \toprule
    Trace ID & Python/reference & BSV/Verilator & Network model & Byte match & Event-order match \\
    \midrule
    To measure & To archive & To archive & To archive & Exact required & Exact required \\
    \bottomrule
  \end{tabular}
\end{table}
}

\subsection{GPU-Ready to On-NIC Payload: A Code-Derived Cycle Budget}
\label{sec:eval:submission-cycles}

The packet-level results above begin after the NIC has accepted a request.  This
subsection isolates the preceding control path because it determines how soon
the first payload read can start.  We statically count the serialized RTL stages
in M0 and M1, then express the external GPU, PCIe, and DMA terms symbolically.
We do not claim a measured GPU or target-FPGA latency.

% Cycle provenance (active tree):
% M0: TxMvp.bsv issueSlotQuery 28--41, completeSlotQuery 43--73,
% issueMrQuery 104--109, completeMrQuery 111--136,
% completeQpReserve 138--162, issuePgtQuery 188--218,
% completePgtQuery 220--258, emitBurst 278--324.
% M1: DmaEngine.bsv acceptFirstChunk 180--223, allocateBurst 303--315,
% installScheduledBurst 357--383, receiveCompletion 398--441,
% drainLastBeat 472--508, stageOrderedBeat 604--670, payloadOut 721--730.
% The active M0 input is already-decoded; GPU fence, credit CAS, MMIO store,
% PCIe transport, and address-to-QPN decode are outside this RTL.

\textbf{Model boundary.}
The active M0 input is \texttt{DecodedGdrpSubmission}; the GPU-side release
that makes payload visible, credit reservation, the 8-B Doorbell store, PCIe
transport, and address-to-QPN decoding are outside the current RTL.  The cycle
count below therefore starts when that decoded submission is visible at the M0
input FIFO and ends when the first \texttt{DmaReadRequest} is accepted at the
M1 interface.  The endpoint for payload arrival is the M1 \texttt{payloadOut}
handshake; compaction, packetization, and M3 are excluded.  A production front
end also needs address/QPN extraction and Doorbell field decode before this
point; we denote that unmodelled term $D_{\mathrm{decode}}$ and do not assign
it a cycle count.

\textbf{Why the M0 path is serialized.}
Each table query supplies the next query's key or permission.  The Slot lookup
returns the local LKEY, remote base, and length limit; the MR lookup validates
that LKEY and returns the PGT offset; the QP reservation returns PMTU and PSN;
the PGT lookup translates the first physical page.  A burst splitter then emits
a 4-KiB-safe, at-most-1-KiB DMA request.  In the functional model each table is
a one-cycle register lookup.  At a target clock, a BRAM-backed PGT may add
cycles; the counts in Table~\ref{tab:eval:submission-stages} are therefore a
lower bound for this RTL organization, not a post-route timing result.

\begin{table}[H]
  \centering
  \caption{Code-derived serialized cycle budget.  The first row starts after a
  decoded submission is visible; all counts assume no stalls and one-cycle
  table lookups.  ``Exported'' includes the M0--M1 handshake.}
  \label{tab:eval:submission-stages}
  \footnotesize
  \begin{tabular}{@{}p{0.66\columnwidth}r@{}}
    \toprule
    Serialized segment & Cycles \\
    \midrule
    M0 Slot, MR, QP, and PGT queries plus DMA split & 12 \\
    M0 adapter and context/chunk synchronizer & 3 \\
    M1 normalizer, scheduler, and reorder installation & 3 \\
    Exported DMA-request handshake & 1 \\
    \midrule
    Total to first external DMA request & 19 \\
    \bottomrule
  \end{tabular}
\end{table}

M0 itself needs 12 cycles: three each for the Slot, MR, and PGT
issue/lookup/join triples, two for the QP reserve/join, and one for the burst
splitter.  Three adapter/synchronizer hops, three M1 installation stages, and
the exported handshake add seven cycles.  This count describes the
\texttt{TxMvp} $\rightarrow$ \texttt{M0MessageAdapter} $\rightarrow$
\texttt{DmaPhaseAIntegration} topology; the standalone Phase-A testbench starts
at its own \texttt{contextIn}/\texttt{chunkIn} and is shorter.  Thus
\begin{equation}
  C_{\mathrm{req}} = 12 + 3 + 3 + 1 = 19
\end{equation}
from submission visibility to the first external DMA-request handshake.  At
250~MHz this is 76~ns; at 500~MHz it is 38~ns.  Both are clock conversions of
the same RTL path, not measured operating frequencies.

\textbf{First payload beat.}
The M1 completion reorder waits until the entire first chunk is received before
it drains any beat: \texttt{recvComplete} becomes true only after the final
128-B completion beat.  If $R$ is the request-to-first-completion delay in
cycles and $K$ is the number of 128-B completion beats in the first DMA chunk,
the first payload beat and the full chunk appear at
\begin{align}
  C_{\mathrm{first}} &= C_{\mathrm{req}} + R + K + 3, \\
  C_{\mathrm{chunk}} &= C_{\mathrm{req}} + R + 2K + 2.
\end{align}
These expressions assume one completion beat per cycle, no downstream
backpressure, and $R$ expressed in the same NIC clock; if the completion
interface uses another clock, convert $R$ to time before adding it.
For a 4-KiB-aligned 1-KiB chunk, $K=8$, so the first beat is $R+30$ and the
complete chunk is $R+37$; an unaligned chunk has $K=9$ and adds one cycle to
each.  The $K+3$ term is a deliberate consequence of whole-chunk reordering,
not a PCIe limit.  A streaming reorder policy could expose the first beat
earlier, but would no longer guarantee that the chunk's beats leave in order.

\textbf{Comparison with IBGDA.}
A conventional GPU-initiated RDMA path first writes a WQE and doorbell record,
then rings the NIC doorbell; the NIC fetches the WQE before it can issue the
payload DMA.  IBGDA and GDAKI move the queue and doorbell into GPU-accessible
memory, but retain WQE construction and NIC WQE fetch~\cite{nvshmem-ibgda,baydamirli2026gpu}.
A recent measurement of minimal GPU-initiated transports reports that
WQE construction, doorbell-record update, and the doorbell store alone cost
0.70~$\mu$s in \emph{mini-gda}, and 5.31~$\mu$s in the NVSHMEM public path on
an H200/CX-7 200-Gb/s InfiniBand platform~\cite{baydamirli2026gpu}.  These are
GPU-side issue times, not NIC-internal cycles.

\begin{table}[H]
  \centering
  \caption{Published p50 issue and put+completion times for one 8-B operation
  on H200/CX-7 200-Gb/s IB~\cite{baydamirli2026gpu}.  kSM-cycle columns are our
  conversion at an assumed 2.0-GHz SM clock, not reported by the authors.}
  \label{tab:eval:ibgda}
  \scriptsize
  \begin{tabular}{@{}lrrrr@{}}
    \toprule
    Path & Issue & Put+comp. & Issue & Put+comp. \\
    & ($\mu$s) & ($\mu$s) & (kSM cyc.) & (kSM cyc.) \\
    \midrule
    mini-gda, inline & 0.70 & 4.03 & 1.4 & 8.1 \\
    GDAKI, inline & 1.60 & 6.02 & 3.2 & 12.1 \\
    NVSHMEM, internal & 4.26 & 9.18 & 8.5 & 18.4 \\
    NVSHMEM, public & 5.31 & 11.10 & 10.6 & 22.2 \\
    \bottomrule
  \end{tabular}
\end{table}

The measured 0.70-$\mu$s issue floor is 1.4k SM cycles at 2~GHz.  The
OpenRDMA M0 path is 12 \emph{NIC} cycles, or 24--60~ns at 200--500~MHz.
These quantities are in different clock domains and must not be subtracted as
cycles.  Their comparison is meaningful only in time: the on-NIC lookup is
small relative to the measured GPU issue floor, so the GPU fence, MMIO write,
and PCIe transport—not the four table lookups—are the likely dominant terms.

The architectural difference is better seen as a transaction count.  For one
1-KiB operation, Table~\ref{tab:eval:transactions} lists the control and data
movement.  IBGDA must publish a WQE and a doorbell record, and the NIC must
fetch the WQE before it can fetch the payload.  OpenRDMA publishes only the
8-B Doorbell; its Slot, MR, QP, and PGT state lives on the NIC, so there is no
WQE fetch.  The NIC-side replacement is the 12-cycle M0 lookup/split path.

\begin{table}[H]
  \centering
  \caption{Submission-path transactions for one 1-KiB operation.  WQE sizes
  are at least one 64-B basic block~\cite{baydamirli2026gpu}; payload is
  fetched by the NIC in both designs.}
  \label{tab:eval:transactions}
  \footnotesize
  \begin{tabular}{@{}p{0.20\columnwidth}p{0.42\columnwidth}p{0.27\columnwidth}@{}}
    \toprule
    Path & GPU-side publication & NIC-side dependent reads \\
    \midrule
    OpenRDMA & 8-B Doorbell & payload only \\
    IBGDA & WQE + dbrec + 8-B Doorbell & WQE, then payload \\
    \bottomrule
  \end{tabular}
\end{table}

Assuming the same GPU release, Doorbell write, PCIe transport, and payload
read, let $T_{\mathrm{IBGDA}}$ and $T_{\mathrm{OpenRDMA}}$ denote data-ready
to first on-NIC payload; the mechanism difference is
\begin{equation}
  T_{\mathrm{IBGDA}} - T_{\mathrm{OpenRDMA}}
  \approx T_{\mathrm{WQE\ build}} + T_{\mathrm{dbrec}}
  + T_{\mathrm{WQE\ fetch}} - \frac{12}{f_{\mathrm{NIC}}}.
\end{equation}
The equation isolates the WQE work that OpenRDMA removes and the local table
path it adds.  It does not predict an end-to-end speedup: the common GPU fence,
MMIO, and PCIe terms are absent from the active RTL and are not measured here.
A same-platform experiment must measure those terms for both designs; comparing
the 0.70-$\mu$s GPU issue floor with NIC RTL cycles directly would be invalid.

The model also explains why the M0 cycle count should not be optimized in
isolation.  At 250~MHz the complete request path is 76~ns, while the measured
minimal GPU issue alone is 700~ns.  The useful architectural claim is therefore
narrow: OpenRDMA removes WQE construction and a dependent WQE fetch from the
NIC's critical path, and replaces them with a 12-cycle on-NIC table path.  It
does not establish a lower GPU-to-NIC latency without a matched
GPU/PCIe/fence measurement.

\subsection{RQ1: Where Does the TX Pipeline Spend Time?}
\label{sec:eval:tx}

Figure~\ref{fig:eval:tx-latency} separates startup from complete message
emission.  At a 4-KiB PMTU, a 64-B request starts its first packet after 102 ns
and finishes after 118 ns.  For 16 KiB, these boundaries are 510 ns and 1,738 ns.
The first-packet delay is also 510 ns at 8 KiB, while latest EOF continues to
increase with message size; startup and drain are therefore distinct costs.

\begin{figure}[tbp]
  \centering
  \includegraphics[width=0.78\textwidth]{\EvalFigureRoot/eval-preview-fig1-tx-latency.pdf}
  \caption{Isolated TX timing from request acceptance in deterministic runs.
  Panel (a) uses a 4-KiB PMTU; panel (b) uses a 2-KiB message.  ``All packets
  emitted'' uses the latest EOF, including when a small tail finishes before an
  earlier large packet.}
  \label{fig:eval:tx-latency}
\end{figure}

Segmentation makes the completion boundary consequential.  In the 4-KiB case,
the 28-B tail reaches EOF at 380 ns, before the 4,068-B first fragment; the
message leaves TX only at 942 ns.  Treating the highest packet index as message
completion would underestimate emission by 562 ns.  For a fixed 2-KiB message,
reducing PMTU from 4,096 B to 256 B changes the packet count from one to nine and
reduces latest EOF from 522 ns to 284 ns in this model.  This result reflects
parallel output lanes; it does not establish a network-optimal PMTU because it
does not price all real wire and receiver costs.

Figure~\ref{fig:eval:tx-rate} next exposes startup/drain and DMA dependence.  For
1-KiB messages, increasing a four-QP batch from 32 to 128 requests raises TX
effective payload rate from 175.2 to 201.9 Gb/s.  The corresponding network
goodputs are 177.9 and 202.8 Gb/s, but their start boundaries differ.  At 64 B,
TX reaches 38.3 and 47.2 Gb/s; at 16 KiB, 123.0 and 124.9 Gb/s.  Thus, larger
messages do not automatically maximize rate in the current pipeline/model.

\begin{figure}[tbp]
  \centering
  \includegraphics[width=0.78\textwidth]{\EvalFigureRoot/eval-preview-fig2-tx-rate.pdf}
  \caption{Finite-batch TX effective payload rate with four active QPs and
  200-ns propagation on each of four paths.  (a) Deterministic 32- and
  128-message batches amortize startup and drain differently.  (b) Added wait
  before each abstract DMA completion lowers the rate of a 32-message, 1-KiB
  workload.}
  \label{fig:eval:tx-rate}
\end{figure}

Adding 40 or 200 ns (20 or 100 model cycles) before each DMA completion lowers
the 1-KiB, 32-message rate from 175.2 Gb/s to 104.1 or 34.3 Gb/s.  Four egress
paths cannot compensate for payload that has not reached packetization.  These
traces establish a DMA-availability dependency, but not its attribution to a
specific RTL queue or physical PCIe transaction.

\EvalPlanned{
\begin{table}[tbp]
  \centering
  \caption{\textbf{P-CYCLE --- PLANNED MEASUREMENT -- NO DATA.}
  TX and RX cycle breakdowns use distinct semantic endpoints.}
  \label{tab:eval:planned-cycle}
  \scriptsize
  \begin{tabular}{@{}lp{0.20\textwidth}lllll@{}}
    \toprule
    TX source & Configuration & Source--DMA req. & DMA done--M2 & M2--commit enqueue & EOF dequeue & Notify delay \\
    \midrule
    M0 or integration & Payload, PMTU, QPs, delay, backpressure & \multicolumn{5}{c}{To measure in model cycles} \\
    \bottomrule
  \end{tabular}

  \vspace{0.4em}
  \begin{tabular}{@{}lp{0.20\textwidth}lllll@{}}
    \toprule
    RX source & Configuration & SOF--parse & Parse--admit & Admit--DMA req. & DMA cpl.--retire & Retire--CQE/Flag \\
    \midrule
    M4 wire ingress & Payload, PMTU, QPs, order, backpressure & \multicolumn{5}{c}{To measure in model cycles} \\
    \bottomrule
  \end{tabular}
\end{table}
}

\subsection{RQ2: When Does Packet Spraying Help?}
\label{sec:eval:spraying}

We compare implemented round-robin spraying among ready paths with fixed
per-QP placement using the low two QPN bits.  Each pair has the same offered
sequence of eight messages, size, QP assignment, 4-KiB PMTU, and four 200-ns
paths.  Policy-dependent backpressure may still change request acceptance and
frame injection.  Fixed placement is therefore an implementation-level
path-selection baseline, not a comparison with a commercial RNIC.

\begin{figure}[tbp]
  \centering
  \includegraphics[width=0.80\textwidth]{\EvalFigureRoot/eval-preview-fig3-path-policy.pdf}
  \caption{Packet spraying versus one fixed path per QP for deterministic
  eight-message batches at a 4-KiB PMTU on four 200-ns paths.  The upper row
  reports batch network goodput; the lower row reports mean
  first-injection-to-last-arrival completion.}
  \label{fig:eval:path-policy}
\end{figure}

With one active QP, fixed placement leaves three paths idle.  Spraying raises
1-KiB goodput from 48.0 to 119.2 Gb/s ($2.48\times$) and 16-KiB goodput from
53.8 to 123.5 Gb/s ($2.30\times$).  For 16 KiB, mean network completion falls
from 2.287 to 1.183 $\mu$s.  A 1-KiB message is one packet whose 287.68-ns
injection-to-arrival latency does not change; its gain comes from placing
successive messages of one QP on different links, not from accelerating a
packet.

As active QPs increase, fixed placement occupies more links.  At four QPs, both
1-KiB policies reach 119.2 Gb/s.  For 4-KiB messages the result reverses: fixed
placement reaches 157.9 Gb/s, while spraying reaches 116.2 Gb/s, 26.4\% lower;
mean completion is 0.538 versus 1.140 $\mu$s.  Each message consists of a
4,068-B fragment and a 28-B tail.  Spraying assigns 12,204, 8,192, 8,164, and
4,208 payload bytes to the four paths, whereas fixed placement assigns 8,192 B
to each.  Packet-opportunity balance is not byte balance, and the small tail
also changes the emission schedule.  The experiment does not isolate these two
causes, so the supported conclusion is conditional: spraying helps when QP
count underutilizes paths, but packet-granular round robin is not uniformly
better near segmentation boundaries.

\subsection{RQ3: How Much Path Skew Can RX Absorb?}
\label{sec:eval:rx-skew}

A 1-KiB message at a 1-KiB PMTU produces 996-B and 28-B packets.  We hold three
paths at 200 ns and add 0--2,000 ns to the fourth.  In a two-message, one-QP
burst, mean network completion is 414.44 ns at 0 and 200 ns of extra delay, then
560.08, 810.08, and 1,310.08 ns at 500, 1,000, and 2,000 ns.  Initially the
delayed tail is not the message's final arrival; once it is, completion follows
that path's additional delay.

\begin{figure}[tbp]
  \centering
  \includegraphics[width=0.78\textwidth]{\EvalFigureRoot/eval-preview-fig4-rx-skew.pdf}
  \caption{Receive behavior for deterministic 1-KiB messages at a 1-KiB PMTU,
  with three 200-ns paths and one slower path.  (a) Mean network completion for
  a two-message, one-QP burst.  (b) Sampled largest passing and smallest failing
  one-QP bursts; markers are not measured occupancy or a general capacity bound.}
  \label{fig:eval:rx-skew}
\end{figure}

Longer bursts expose a separate receiver limit.  At equal path delay, 11
messages pass and 12 fail; with 1,000 or 2,000 ns of extra delay, nine pass and
ten fail.  At 200 and 500 ns the sampled pairs are 13/14 and 12/13.  Failure is
\texttt{RX\_OUT\_OF\_WINDOW}, not network loss: every frame arrives in the
recorded failing runs.  Equal propagation also does not guarantee in-order RTL
input because small tails and packets on different lanes have different
serialization histories.  The non-monotonic 200-ns point is a property of this
schedule, not evidence that delay increases capacity.  Without a closed
retransmission loop, a rejected packet cannot be recovered in this batch.

\EvalPlanned{
\begin{table}[tbp]
  \centering
  \caption{\textbf{P-SCALE --- PLANNED MEASUREMENT -- NO DATA.}
  Steady-state and robustness sweep; stochastic rows require fixed seed sets
  and 95\% confidence intervals.}
  \label{tab:eval:planned-scale}
  \footnotesize
  \begin{tabular}{llllllll}
    \toprule
    QPs & Window & Burst/loss/skew & Warm-up & Measure & Drain & Goodput/FCT & Occupancy/stall/failure \\
    \midrule
    To sweep & To sweep & To sweep & To define & To define & Excluded & P50/P95/P99 & To record \\
    \bottomrule
  \end{tabular}
\end{table}
}

\subsection{RQ4: What Does Ordered Completion Cost?}
\label{sec:eval:ordering}

We isolate the Ordered Queue because packet arrival alone does not authorize an
Ordered Flag.  The testbench enqueues a Flag, injects retirement evidence at a
controlled barrier, and optionally stalls the output.  The retirement event is
an RTL input in this experiment, not a measured GPU-memory DMA completion.

\begin{figure}[tbp]
  \centering
  \includegraphics[width=0.78\textwidth]{\EvalFigureRoot/eval-preview-fig5-ordered-completion.pdf}
  \caption{Controlled ordering tests.  (a) Flag release with delayed retirement
  and output backpressure; one model cycle is 2 ns.  (b) Four QPs complete in
  order 3,2,1,0 but retire in order 0,1,2,3; delaying QP 0 transfers its delay to
  already-completed QP 3.}
  \label{fig:eval:ordered-completion}
\end{figure}

With no output stall, retirement delays of 0, 20, and 100 cycles produce observed
Flag enqueue-to-acceptance times of 10, 50, and 210 ns.  Combining 20 cycles of
retirement delay with 20 cycles of output stall yields 90 ns.  In the 100-cycle
case, retirement evidence is accepted at 202 ns, the request becomes ready at
208 ns, and it is accepted at 210 ns; credit is dequeued at 212 ns and the
modeled response at 214 ns.  The two-nanosecond gaps include sequential
testbench drive/drain cadence and are not intrinsic interface-latency claims.
Capacity tests return one, two, and four credits for the corresponding released
Flags and detect a fifth insertion into a blocked four-entry queue.

The second experiment places one request from each of four QPs on different DMA
lanes and completes them in reverse order.  Retirement remains 0,1,2,3.  QP 3,
which completes first, waits 34 ns in the baseline and 1,034 ns when QP 0 is
delayed by 1,000 ns.  QP 0's own completion-to-retire interval remains 6 ns.
Thus a delayed global head transfers essentially all of its delay to unrelated,
already-completed QPs.  Per-QP retirement may remove this cross-QP HoL, but its
state and arbitration cost have not been evaluated.

\subsection{RQ5: Established Cost and Remaining Evidence}
\label{sec:eval:cost}

Table~\ref{tab:eval:logical-state} groups the declared entry data in the
archived RTL inventory.  The listed structures contain 2,586,432 logical bits
(315.73 KiB).  RX payload-frame storage contributes 260 KiB; the four per-QP
Ordered Queues together hold only 416 B of entry data.  Within this inventory,
payload buffering rather than Flag metadata dominates.  Increasing Flag queue
space would not address receive-window pressure or global-retirement blocking.

\begin{table}[tbp]
  \centering
  \caption{Logical state footprint of listed structures.  Counts are declared
  entry data, not placed FPGA memory or a complete-tree resource estimate.}
  \label{tab:eval:logical-state}
  \footnotesize
  \begin{tabular}{p{0.61\columnwidth}rr}
    \toprule
    Structure & Logical bits & KiB \\
    \midrule
    RX payload frame storage & 2,129,920 & 260.00 \\
    TX channel payload storage & 262,144 & 32.00 \\
    TX DMA completion reorder & 158,848 & 19.39 \\
    RX pending-packet table & 13,888 & 1.70 \\
    Per-QP Ordered Queues & 3,328 & 0.41 \\
    Other listed state & 18,304 & 2.23 \\
    \midrule
    Total listed capacity & 2,586,432 & 315.73 \\
    \bottomrule
  \end{tabular}
\end{table}

Logical bits are not FPGA resource utilization: synthesis may repack memories
and add FIFOs, control, routing, and vendor-IP cost.  We therefore keep this
inventory separate from target-device implementation evidence, as is necessary
in FPGA system evaluations~\cite{wang2023cxloverethernet}.  A historical
standalone M1 synthesis artifact is excluded because its Fitter failed and it
has no placed timing result.  Credible FCCM implementation evidence still
requires a controlled top and SDC, successful synthesis/Fitter/TimeQuest,
module and total ALM/LUT/FF/BRAM/DSP usage, target clock, WNS/TNS, and post-route
Fmax.  Target-clock feasibility must not be inferred from the 2-ns model cycle.

\EvalPlanned{
\begin{table}[tbp]
  \centering
  \caption{\textbf{P-SYNTH --- PLANNED MEASUREMENT -- NO DATA.}
  Existing module/tops are reported separately; a full M0--M11 top is blocked
  because the three integrations are not connected.}
  \label{tab:eval:planned-synth}
  \scriptsize
  \begin{tabular}{@{}lllllll@{}}
    \toprule
    Component/top & Evidence stage & Status & ALM/LUT--FF & RAM--DSP & Target--WNS/TNS & Fmax \\
    \midrule
    M0, M1, M2, M3 & Post-route required & Not run & -- & -- & -- & -- \\
    M4--M8, \texttt{RxTop} & Post-route required & Not run & -- & -- & -- & -- \\
    M9--M11, \texttt{RetxTop} & Post-route required & Not run & -- & -- & -- & -- \\
    Full M0--M11 top & N/A & Blocked: top does not exist & -- & -- & -- & -- \\
    \bottomrule
  \end{tabular}
\end{table}
}

\begin{table}[tbp]
  \centering
  \caption{Evidence closure for an FCCM implementation evaluation.  The right
  column states required measurements, not inferred results.}
  \label{tab:eval:evidence-closure}
  \scriptsize
  \begin{tabular}{@{}p{0.25\columnwidth}p{0.31\columnwidth}p{0.34\columnwidth}@{}}
    \toprule
    Dimension & Evidence here & Required closure \\
    \midrule
    Function & Archived co-simulation and controlled RTL tests & Restore trace-level cross-model artifact \\
    Cycle cost & Endpoint intervals under abstract DMA & Stage, stall, and occupancy counters \\
    Scale/robustness & Finite batches and sampled skew boundary & Warm-up-separated, multi-seed QP/loss sweeps \\
    Resources & 315.73 KiB listed logical state & Successful Fitter/TimeQuest and post-route report \\
    Recovery & M9--M11 functional seam outside this batch & Connect M3, M6, runtime, and retry egress \\
    Physical system & Not measured & PCIe/MAC/GPU throughput, latency, and power \\
    \bottomrule
  \end{tabular}
\end{table}

The present evidence establishes functional datapath execution, conditional
multipath gains, schedule-specific sampled receive-window pass/fail boundaries,
cross-QP ordering cost, and a logical-state baseline.  It does not establish physical-NIC throughput,
GPU-visible completion latency, closed-loop loss recovery, SRv6 cost,
target-FPGA frequency, measured power, or application speedup.  Completing the
evaluation requires (i) event-level closure among the reference model, RTL, and
restored network harness; (ii) explicit pipeline latency, occupancy, and stall
counters; (iii) warm-up/steady-state/drain-separated QP, window, skew, and loss
sweeps; and (iv) successful target-device place-and-route followed by physical
PCIe/MAC/GPU measurements.  These additions strengthen the evidence level
without changing the negative results and scope boundaries reported above.
